\subsubsection{Descente des opérations au quotient d’évaluation}
\label{sec:ch-descenso-fibras}

Le quotient précédent identifie les historiques exprimant la même valeur, et non leurs registres complets. Le critère suivant caractérise exactement quand une opération sur les historiques induit une opération sur ce quotient sans choisir d'historique privilégié dans chaque fibre ; sa formulation prolonge la construction généalogique établie dans les sections précédentes.

\begin{proposition}[Critère de descente par fibres d’évaluation]
\label{prop:ch-descenso-fibras}
Soit \(q:\Omega\to J\) l'évaluation sur son image d'une coordinatisation généalogique. Soit \(\star:D\to\Omega\) une opération avec \(D\subseteq\Omega^2\) saturé par les fibres de \(q\times q\). Il existe une unique opération \(\overline\star:(q\times q)(D)\to J\) telle que
\[
 q(x\star y)=q(x)\,\overline\star\,q(y)
\]
si et seulement si, pour des paires du domaine,
\[
 q(x)=q(x'),\quad q(y)=q(y')
 \quad\Longrightarrow\quad q(x\star y)=q(x'\star y').
\]
\end{proposition}

\begin{proof}
Si l'opération visible existe, son résultat dépend seulement des deux valeurs, ce qui prouve la nécessité. Pour la suffisance, on définit \(s\,\overline\star\,t=q(x\star y)\), où \(q(x)=s\) et \(q(y)=t\). La saturation conserve le domaine lors du changement de représentants ; la condition énoncée conserve le résultat. La surjectivité de \(q\times q\) sur l'image du domaine prouve l'unicité.
\end{proof}

Une coordonnée ne remplace pas automatiquement un historique comme argument de toutes ses opérations : le critère détermine quand cette substitution est légitime. En particulier, transporter une opération à travers une bijection ne suffit pas à l'identifier à une opération native antérieure à la lecture. Si une coordinatisation \(\Omega\) est compacte et si \(q\) est continue et surjective sur un espace séparé \(J\), l'évaluation induit en outre un homéomorphisme \(\Omega/{\sim_q}\to J\) : c'est une bijection continue d'un compact vers un espace séparé. C'est le \emph{quotient}, et pas nécessairement l'espace des historiques, qui est identifié topologiquement à \(J\).
